\documentclass[11pt]{article}
\usepackage[a4paper,margin=1in]{geometry}
\usepackage{amsmath,amssymb,amsthm,mathtools}

\newtheorem{theorem}{Theorem}[section]
\newtheorem{lemma}[theorem]{Lemma}
\newtheorem{proposition}[theorem]{Proposition}
\newtheorem{conjecture}[theorem]{Conjecture}
\newtheorem{remark}[theorem]{Remark}

\newcommand{\Om}{\Omega}
\newcommand{\R}{\mathbb R}
\newcommand{\Sph}{\mathbb S}
\newcommand{\cA}{\mathcal A}
\newcommand{\dT}{\delta_T}

\title{Hybrid Level-Flow Route to an \(L^\infty\) Graph}
\author{}
\date{May 18, 2026}

\begin{document}
\maketitle

\section{The Intended Hybrid}

The intended route is not a pure Plan 3 single-level argument and not the
full Plan 2 foliation argument.  It is the following hybrid.

\begin{enumerate}
\item Move inward along torsion levels \(E_t=\{u_\Om>t\}\).
\item Stop at a level \(E_{\hat t}\) for which the discarded layer is still
sharp:
\[
   |\Om\setminus E_{\hat t}|\le C_n\dT(\Om).
\]
\item Prove that this level is an \(L^\infty\) radial graph after translation
and volume normalization:
\[
   E_{\hat t}^{\sharp}
      =\{z+r\theta:0<r<1+\varphi(\theta)\},
   \qquad
   \|\varphi\|_{L^\infty(\Sph^1)}\le\eta_0.
\]
No \(C^1\), curvature, perimeter, or \(H^1\) control is required.
\item Apply the direct small-amplitude radial-graph theorem:
\[
   \dT(E_{\hat t})\ge c\|\varphi\|_{L^2}^2
      \ge c'\cA(E_{\hat t})^2 .
\]
\item Transfer back using the superlevel set estimate
\[
   \cA(\Om)\le \cA(E_{\hat t})+
      2|\Om\setminus E_{\hat t}|/|\Om|.
\]
\end{enumerate}

The point of the direct radial-graph theorem is exactly that Step 3 only has
to produce a small-amplitude graph.  The usual high-norm graph entry is not
needed.

\section{The Conditional Assembly}

The following theorem is already enough to close the hybrid in dimension
\(2\).

\begin{conjecture}[Level-flow \(L^\infty\) graph-entry theorem in \(n=2\)]
\label{conj:linfty-entry}
There are computable constants \(\delta_0,\eta_0,C>0\) such that, if
\(\Om\subset\R^2\), \(|\Om|=\pi\), and \(\dT(\Om)\le\delta_0\), then there
exists a regular value \(\hat t\) of \(u_\Om\) and a center \(z\in\R^2\) such
that
\[
   |\Om\setminus E_{\hat t}|\le C\dT(\Om),
\]
and, after translating by \(z\) and rescaling to area \(\pi\),
\[
   E_{\hat t}^{\sharp}
      =\{r e^{i\theta}:0<r<1+\varphi(\theta)\},
   \qquad
   \|\varphi\|_{L^\infty(\Sph^1)}\le\eta_0,
\]
with the volume and first moment normalized.
\end{conjecture}

\begin{proposition}[Conditional closure of the hybrid route]
Assume Conjecture~\ref{conj:linfty-entry}.  Then the sharp quantitative
Saint--Venant estimate
\[
   \cA(\Om)^2\le C_2\,\dT(\Om)
\]
holds in dimension \(2\), with computable constants.
\end{proposition}

\begin{proof}
Choose \(E_{\hat t}\) from Conjecture~\ref{conj:linfty-entry}.  Step 3 of
Plan 3 gives the exact deficit transfer
\[
   \dT(E_{\hat t})\le
     \frac{1}{1-|\Om\setminus E_{\hat t}|/|\Om|}\,\dT(\Om)
     \le 2\dT(\Om)
\]
after reducing \(\delta_0\).  The direct radial-graph theorem applies to
the normalized graph \(E_{\hat t}^{\sharp}\) and gives
\[
   \cA(E_{\hat t})^2=\cA(E_{\hat t}^{\sharp})^2
      \le C\,\dT(E_{\hat t}^{\sharp})
      = C\,\dT(E_{\hat t})
      \le C\,\dT(\Om),
\]
where the equality uses scale invariance of the normalized deficit.  Finally,
Step 5 gives
\[
   \cA(\Om)\le \cA(E_{\hat t})+
        2|\Om\setminus E_{\hat t}|/|\Om|
     \le C\dT(\Om)^{1/2}+C\dT(\Om).
\]
Squaring and reducing \(\delta_0\) gives the claim.
\end{proof}

\section{What Is Already Proved}

\begin{itemize}
\item The \(O(\dT)\) layer selection is proved in Plan 3 Step 1, using the
volume variable \(s=m(t)\).  This avoids the old \(t\)-length gap.
\item The deficit transfer to \(E_{\hat t}\) is proved in Plan 3 Step 3,
using the exact identity \(u_\Om-\hat t=u_{E_{\hat t}}\).
\item The transfer back to \(\Om\) is proved in Plan 3 Step 5.
\item The small-amplitude radial-graph theorem is proved in
\texttt{H12\_ASYMMETRY\_DIRECT\_PROOF.tex}.  It needs only
\(\|\varphi\|_{L^\infty}\) small and \(\|\varphi\|_{L^2}\) small; it does
not require a slope or curvature bound.
\end{itemize}

\section{What the Entry Theorem Must Prove}

Conjecture~\ref{conj:linfty-entry} is stronger than mere annular
containment.  It must prove both:
\[
   B_{1-\eta_0}(z)\subset E_{\hat t}^{\sharp}
      \subset B_{1+\eta_0}(z),
\]
and one radial crossing from \(z\) in every direction.  These are separate
issues.

\subsection{Outward fingers}

The ball-plus-tube model supports the plan in dimension \(2\).  For a tube
of width \(w\) and length \(L=O(1)\),
\[
   \dT(\Om)\simeq wL,
\]
while clearing the tube by passing above its cap height costs a ball-shell
volume of order \(w^2\), plus the tube volume \(wL\).  Since \(w^2\le CwL\)
in the nontrivial regime, the tube can be cleared within an \(O(\dT)\)
discarded layer in \(n=2\).  Also, on active tube levels,
\[
   D_H(t)\simeq L/w,
\]
so those levels are very expensive for the \(D_H\)-identity.  This is the
main reason the \(n=2\) hybrid remains plausible.

\subsection{Inward fjords and holes}

The harder part is inward geometry.  Since \(E_t\subset\Om\), moving inward
cannot fill a missing fjord or hole.  Therefore the entry theorem must show
that a small-deficit planar domain cannot have a fjord penetrating a fixed
distance into the ball, unless that fjord is already paid for by
\(\dT(\Om)\) at a fixed scale.

The natural tool is planar capacity/modulus.  A narrow slit of fixed depth
has positive \(W^{1,2}\)-capacity in the disk, so imposing \(u=0\) on such a
slit should lower torsional rigidity by a definite amount, independently of
the slit width.  A quantitative version would imply:

\[
   \hbox{inward penetration depth}\ge \eta
      \quad\Longrightarrow\quad
   \dT(\Om)\ge c(\eta)>0 .
\]

This would give the inner-radius half of the \(L^\infty\) graph entry after
choosing the small-deficit threshold below \(c(\eta_0)\).  This capacity
lemma is not currently proved in the repo.

\subsection{Single crossing}

Annular containment still permits multiple crossings on a ray.  To feed the
radial-graph theorem one must rule out those crossings, or replace the graph
theorem by a general annular-set theorem.

For the hybrid as stated, the required lemma is:

\[
   \hbox{multiple crossings before the \(O(\dT)\) stop}
      \quad\Longrightarrow\quad
   \hbox{either large }D_H\hbox{ on a positive level interval}
   \hbox{ or fixed torsion loss}.
\]

Plan 2's good-ray/bad-ray slicing machinery is relevant here: bad rays are
measured by multiplicity and can be absorbed by perimeter-type excesses.  But
the current repo does not contain the torsion-energy lower bound that turns
this multiplicity into the \(L^\infty\) graph-entry conclusion.

\section{What Not To Use}

The hybrid should not use the single-level Hessian Chain B.  The false step is
\[
   \dT(D)\gtrsim \int_D |D^2u_D+\tfrac1nI|^2,
\]
which fails on high-frequency smooth near-balls by the spectral calculation
recorded in \texttt{PLAN3\_N2\_SINGLE\_LEVEL\_RECHECK.tex}.  The correct
scale is the radial-graph theorem's \(L^2\)/Fraenkel scale, not an \(H^1\)
Hessian scale.

\section{Precise Status}

The hybrid route is coherent and sharper than the previous annular-set
formulation: if the stopped level is truly an \(L^\infty\) radial graph, the
existing radial-graph theorem finishes the proof.  The missing theorem is
exactly Conjecture~\ref{conj:linfty-entry}.

The proof of that conjecture should likely split into:
\begin{enumerate}
\item a \(D_H\)-integral clearing lemma for outward fingers, using the
\(n=2\) scaling \(D_H\simeq L/w\) on active tube levels;
\item a planar capacity/modulus exclusion of inward fjords of fixed depth;
\item a same-center bad-ray/multiplicity argument to promote annular
containment to one radial crossing.
\end{enumerate}

None of these three ingredients is currently written as a proof in the repo.
The model computations support them in \(n=2\), but they are the actual work
left by the hybrid plan.

\end{document}
